%!TEX TS-program = xelatex
%!TEX root = ../Topo2_WS15/topologie_2.tex
\RequirePackage{fix-cm} 
\documentclass[a4paper, twoside, headsepline, index=totoc,toc=listof, fontsize=10pt, cleardoublepage=empty, headinclude, DIV=12, BCOR=5mm, titlepage,draft]{scrartcl}
\usepackage{scrtime} % KOMA, Uhrzeit ermoeglicht

\usepackage{etoolbox}
\usepackage{letltxmacro}
\usepackage{ifthen}

%--Farbdefinitionen
%-- muss vor tikz geladen werden
\usepackage[usenames, table, x11names]{xcolor}
\definecolor{dark_gray}{gray}{0.45}
\definecolor{light_gray}{gray}{0.6}
\usepackage[final]{graphicx}

%--Zum Zeichnen
%-- muss vor polyglossia bzw. babel geladen werden
\usepackage{tikz}
\usepackage{tikz-cd}
\usetikzlibrary{external}
\tikzset{>=latex}
\usetikzlibrary{shapes,arrows.meta,intersections}
\usetikzlibrary{calc,3d}
\usetikzlibrary{decorations.pathreplacing,decorations.markings, decorations.pathmorphing}
\usetikzlibrary{angles}

%-- Konfiguration von tikzexternalize
\tikzexternalize[prefix=tikz/,up to date check=diff]
\pgfkeys{/pgf/images/include external/.code=\includegraphics{#1}}
\tikzset{external/system call={xelatex \tikzexternalcheckshellescape -halt-on-error -interaction=batchmode --shell-escape -jobname "\image" "\texsource"}}


%-- tikzexternalize fuer tikzcd deaktivieren, da inkompatibel
\AtBeginEnvironment{tikzcd}{\tikzexternaldisable}
\AtEndEnvironment{tikzcd}{\tikzexternalenable}

%-- um Inkompatibilitaeten von quotes und polyglossia bzw. babel zu vermeiden
\tikzset{
  every picture/.append style={
    execute at begin picture={\shorthandoff{"}},
    execute at end picture={\shorthandon{"}}
  }
}
\usetikzlibrary{quotes}
\usepackage{pgfplots}
\usepgfplotslibrary{colormaps}
% \pgfplotsset{compat=1.12}



%-- Mathesymbole etc.
\usepackage{mathtools} % beinhaltet amsmath
\mathtoolsset{showonlyrefs, centercolon,showmanualtags}
\newtagform{brackets}[\textbf]{[}{]}
\usetagform{brackets}
\usepackage{fix-cm}
% \usepackage{wasysym} % zusätzliche Symbole
\usepackage{amssymb} % zusätzliche Symbole
% \usepackage{stmaryrd} % für Blitz
\usepackage{nicefrac} % schräge Brüche
\usepackage{faktor}
\newcommand{\Faktor}[1]{\faktor[\textstyle]{#1}}
\usepackage{xfrac}
\usepackage{cancel}
\usepackage{mathdots} % Verbesserung von Punkten wie zB \ldots
\usepackage[bb=px]{mathalfa}

%-- einzelne Symbole aus dem mathx-package
% \DeclareFontFamily{U}{mathx}{\hyphenchar\font45}
% \DeclareFontShape{U}{mathx}{m}{n}{<-> mathx10}{}
% \DeclareSymbolFont{mathx}{U}{mathx}{m}{n}
% \DeclareMathSymbol{\bigplus}{1}{mathx}{"90}

%-- 
\usepackage{silence}
\WarningFilter{latexfont}{Size substitutions with differences}
\WarningFilter{latexfont}{Font shape `U/bbold/m/n' in size}

\DeclareSymbolFont{bbold}{U}{bbold}{m}{n}
\DeclareSymbolFontAlphabet{\mathbbold}{bbold}
\newcommand{\ind}{\mathbbold{1}} % charakteristische-Funktion-Eins

\def\mathul#1#2{\color{#1}\underline{{\color{black}#2}}\color{black}} %farbiges Untersteichen im Mathe-Modus
\renewcommand{\le}{\leqslant}
\renewcommand{\ge}{\geqslant}

%-- Underbrace als Befehl in LaTeX-Syntax (und ohne Spacingprobleme mit nachfolgenden Operatoren...)
\newcommand{\Underbrace}[2]{{\underbrace{#1}_{#2}}}
\newcommand{\Underbracket}[2]{{\underbracket[0.7pt][2pt]{#1}_{#2}}}

%-- Alles was mit Schrift und XeTeX zu tun hat
\usepackage{palatino}
\usepackage[euler-digits]{eulervm}
\usepackage[no-math]{fontspec}
\usepackage{polyglossia} % moderner babel-ersatz
\setmainlanguage[spelling=new,babelshorthands=true]{german}
\shorthandoff{"}
\setotherlanguage{english}
\defaultfontfeatures{Mapping=tex-text, WordSpace={1.2}, Ligatures={Required,Common,Contextual}} %

%-- Alle hier aufgeführten Schriftarten sind in jeder TeXlive-Installation vorhanden!!
%-- Hauptschriftart festlegen
% \setmainfont{texgyrepagella}[Extension=.otf,UprightFont=*-regular,BoldFont=*-bold,ItalicFont=*-italic,BoldItalicFont=*-bolditalic,ItalicFeatures={Style=Historic}]
\setmainfont{TeXGyrePagellaX}[Extension=.otf,UprightFont=*-Regular,BoldFont=*-Bold,ItalicFont=*-Italic,BoldItalicFont=*-BoldItalic,ItalicFeatures={Style=Historic},Ligatures={Required,Common,Contextual,Historic}]
% \setmainfont[BoldFont={* Bold}, ItalicFont={* Light Italic},SmallCapsFont={Linux Libertine O}, SmallCapsFeatures={Letters=SmallCaps}]{Source Sans Pro}


%-- Serifenlose Schriftart festlegen
% \setsansfont{Roboto}[Scale=MatchUppercase,Extension=.ttf,UprightFont=*-Regular,BoldFont=*-Bold,ItalicFont=*-RegularItalic,BoldItalicFont=*-MediumItalic]
% \setsansfont{RobotoSlab}[Scale=MatchUppercase,Extension=.ttf,UprightFont=*-Regular,BoldFont=*-Bold,ItalicFont=*-Thin,SmallCapsFont=*-Regular]
% \setsansfont{FiraSans}[Scale=MatchUppercase,Extension=.otf, UprightFont=*-Regular, BoldFont=*-SemiBold, ItalicFont=*-RegularItalic, BoldItalicFont=*-MediumItalic]
% \setsansfont{SourceSansPro}[Scale=MatchUppercase,Extension=.otf,UprightFont=*-Regular,BoldFont=*-Semibold,ItalicFont=*-LightIt,BoldItalicFont=*-SemiboldIt]
% \setsansfont{texgyreadventor}[Scale=MatchUppercase, Extension=.otf, UprightFont=*-regular, BoldFont=*-bold, ItalicFont=*-italic,BoldItalicFont=*-bolditalic]
% \setsansfont{Kurier}[Scale=MatchUppercase, Extension=.otf, UprightFont=*Medium-Regular, BoldFont=*Heavy-Regular, ItalicFont=*Medium-Italic,BoldItalicFont=*Heavy-Italic]
% \setsansfont{Iwona}[Scale=MatchUppercase, Extension=.otf, UprightFont=*Medium-Regular, BoldFont=*Heavy-Regular, ItalicFont=*Medium-Italic,BoldItalicFont=*Heavy-Italic]
% \setsansfont{LinBiolinum}[Scale=MatchUppercase, Extension=.otf, UprightFont=*_R, BoldFont=*_RB, ItalicFont=*_RI,BoldItalicFont=*_RBO]
% \setsansfont{texgyrebonum}[Scale=MatchUppercase,Extension=.otf, UprightFont=*-regular, BoldFont=*-bold, ItalicFont=*-italic, BoldItalicFont=*-bolditalic]
\setsansfont{texgyreadventor}[Scale=MatchUppercase,Extension=.otf, UprightFont=*-regular, BoldFont=*-bold, ItalicFont=*-italic, BoldItalicFont=*-bolditalic]



% \newfontfamily{\geometric}{texgyreheros}[Scale=MatchLowercase, Extension=.otf, UprightFont=*-regular, BoldFont=*-bold, ItalicFont=*-italic]
% \addtokomafont{disposition}{\geometric}
\setmonofont{SourceCodePro}[Scale=0.9,Extension=.otf,UprightFont=*-Regular, BoldFont=*-Semibold, ItalicFont=*-Light]
\usepackage{xltxtra}
\usepackage{fontawesome}
\usepackage[final]{microtype}


\usepackage[draft=false]{scrlayer-scrpage} % wie fancyhdr, nur optimiert auf KOMA-Skript, leicht andere Syntax

\flushbottom


%--Mixed
\usepackage[shortlabels,inline]{enumitem}
\setlist[itemize,1]{label=\faCaretRight}
\setlist[enumerate]{font=\bfseries}
\setlist[description]{font=\normalfont\bfseries}
\usepackage[german=quotes]{csquotes}
\usepackage{makeidx}
\usepackage{wrapfig}
\usepackage{float}
\usepackage[margin=10pt, font=small, labelfont={sf, bf}, format=plain, indention=1em]{caption}
\captionsetup[wrapfigure]{name=Abb. }
\usepackage{stackrel}


%--Unterstreichung
\usepackage[normalem]{ulem}
\setlength{\ULdepth}{1.8pt}

%--Indexverarbeitung
\newcommand{\bet}[1]{\emph{#1}}
\newcommand{\Index}[1]{\bet{#1}\index{#1}}
\makeindex
\setindexpreamble{{\noindent\sffamily\small Die \emph{Seitenzahlen} sind mit Hyperlinks versehen und somit anklickbar}  \par \bigskip}
\renewcommand{\indexpagestyle}{scrheadings}


%--Marginnote & todonotes
\deffootnote[1.5em]{1.5em}{1.5em}{\textsuperscript{\thefootnotemark}\ }
\usepackage[fulladjust]{marginnote}
\renewcommand*{\marginfont}{\color{dark_gray} \itshape \footnotesize}
\usepackage[textsize=small]{todonotes}
\usepackage{ragged2e}
\renewcommand*{\raggedleftmarginnote}{\RaggedLeft}
\renewcommand*{\raggedrightmarginnote}{\RaggedRight}

%--Todonotes mit tikz/externalize kompatibel machen
\LetLtxMacro{\oldtodo}{\todo}
\renewcommand{\todo}[2][]{\tikzexternaldisable\oldtodo[#1]{#2}\tikzexternalenable}
\LetLtxMacro{\oldmissingfigure}{\missingfigure}
\renewcommand{\missingfigure}[2][]{\tikzexternaldisable\oldmissingfigure[{#1}]{#2}\tikzexternalenable}

%--Konfiguration von Hyperref pdfstartview=FitH, 
\usepackage[hidelinks, pdfpagelabels,  bookmarksopen=true, bookmarksnumbered=true, linkcolor=black, urlcolor=SkyBlue2, plainpages=false,pagebackref, citecolor=black, hypertexnames=true, pdfauthor={Jannes Bantje}, pdfborderstyle={/S/U}, linkbordercolor=SkyBlue2, colorlinks=false,final]{hyperref}


\newcommand{\appendLink}[1]{#1\,\faExternalLink}
\newcommand{\hrefsym}[2]{\href{#1}{\texttt{\appendLink{#2}}}}
\newcommand{\hrefsymmail}[2]{\href{#1}{\texttt{\faEnvelopeO\,#2}}}
\renewcommand{\url}[1]{\hrefsym{#1}{\nolinkurl{#1}}}



% -- QR-Codes
\usepackage{qrcode} %-- hinter hyperref laden!

%--Römische Zahlen
\newcommand{\RM}[1]{\MakeUppercase{\romannumeral #1{}}}

%%--Abkürzungen etc.
\newcommand{\light}{\text{\Large $\lightning$}}

%-- Definitionen von weiteren Mathe-Befehlen

\DeclareMathOperator{\im}{im} %Bild
\DeclareMathOperator{\id}{id} %identische Abbildung
\DeclareMathOperator{\conj}{conj} %Konjugation
\DeclareMathOperator{\sgn}{sgn} %Signum
\DeclareMathOperator{\End}{End} %Endomorphismen
\DeclareMathOperator{\Hom}{Hom} %Homomorphismen
\DeclareMathOperator{\Iso}{Iso} %Isomorphismen
\DeclareMathOperator{\Span}{span} %Span

\DeclareMathOperator{\grad}{grad} %Gradient
\DeclareMathOperator{\dive}{div} %Gradient
\DeclareMathOperator{\rot}{rot} %Rotation


\DeclareMathOperator{\Char}{char} %Charakteristik
\DeclareMathOperator{\Aut}{Aut} %Automorphismen

% -- Zum Finetuning von Befehlen
\makeatletter
\newcommand{\raisemath}[1]{\mathpalette{\raisem@th{#1}}}
\newcommand{\raisem@th}[3]{\raisebox{#1}{$#2#3$}}
\makeatother

\newcommand{\sing}{{\raisemath{1.1pt}{\scriptscriptstyle\mathrm{sing}}}}
\newcommand{\pt}{\mathrm{pt}}
\newcommand{\op}{\mathrm{op}}
\DeclareMathOperator{\Sp}{Sp}
\DeclareMathOperator{\Keg}{Keg}
\newcommand{\slashedi}{i\hspace{-3.5pt}/}
\newcommand{\cupp}{\smallsmile}
\newcommand{\capp}{\smallfrown}
\DeclareMathOperator*{\colim}{colim}


\newcommand{\D}{\ensuremath{\mathrm{D}\mkern-1.5mu}}
\newcommand{\mathd}{\ensuremath{\mathrm{d}\mkern-1.0mu}}
\newcommand{\Tmap}{\ensuremath{\mathrm{T}\mkern-0.85mu}}
\newcommand{\opL}{\ensuremath{\mathrm{L}\mkern-0.6mu}}


\DeclareMathOperator{\supp}{supp} %Träger


% -- Befehle für FunkAna und Operatoralgebren
\newcommand{\tr}{\mathrm{tr}}
\newcommand{\w}{\mathrm{w}}
\newcommand{\sa}{\mathrm{sa}}
\newcommand{\so}{\mathrm{s.o.}}
\newcommand{\weakT}[1]{\ensuremath{\mathcal{T}_{#1}^{\mkern+1.0mu\text{\raisebox{0.4ex}{$\mathrm{w}$}}}}}
\newcommand{\weakTstar}[1]{\ensuremath{\mathcal{T}_{#1}^{\mkern+1.0mu\text{\raisebox{0.4ex}{$\mathrm{w}$}}^*}}}
\newcommand{\finSub}{\subset\mkern-0.7mu \subset}
\DeclareMathOperator{\Inv}{Inv}
\DeclareMathOperator{\dist}{dist}
\newcommand{\simm}{{\hspace{-1.6pt}\raisemath{0.5pt}{\sim}}}
\DeclareMathOperator{\ev}{ev}
\DeclareMathOperator{\Alg}{Alg}


% -- Kategorien
\newcommand{\TOP}{\textsc{Top}}
\newcommand{\HTOP}{\textsc{HTop}}
\newcommand{\VR}{\textsc{VR}}
\newcommand{\MOD}{\textsc{Mod}}
\newcommand{\MONOIDE}{\textsc{Monoide}}
\newcommand{\MENGEN}{\textsc{Mengen}}
\newcommand{\MAN}{\textsc{Man}}
\newcommand{\GRUPPEN}{\textsc{Gruppen}}
\newcommand{\ABELGRUPPEN}{\textsc{Abel.Gruppen}}
\newcommand{\KAT}{\textsc{Kat}}
\newcommand{\FUN}{\textsc{Fun}}
\newcommand{\SIMP}{\textsc{Simp}}
\newcommand{\VEKT}{\textsc{Vekt}}

\DeclareMathOperator{\cell}{cell}
\DeclarePairedDelimiter{\homologieklasse}{\llbracket}{\rrbracket}
\newcommand{\rand}[1]{\ensuremath{\partial^{\scriptscriptstyle #1}}}


% -- theorem packages
\usepackage{amsthm}
\usepackage{thmtools,thm-restate}
\renewcommand{\listtheoremname}{Übersicht aller Aussagen}

% -- Theoreme als PDF-Lesezeichen
\usepackage{bookmark}
\bookmarksetup{open,numbered}
\makeatletter
\newcommand*{\theorembookmark}{%
  \bookmark[
    dest=\@currentHref,
    rellevel=1,
    keeplevel,
  ]{%
    \thmt@thmname\space\csname the\thmt@envname\endcsname
    \ifx\thmt@shortoptarg\@empty
    \else
      \space(\thmt@shortoptarg)%
    \fi
  }%
}   
\makeatother

% -- Definition der einzelnen Umgebungen
\declaretheoremstyle[%
	headfont=\sffamily\bfseries,
	notefont=\normalfont\sffamily\scshape,
	bodyfont=\normalfont,
	headformat=\NUMBER\ \NAME\NOTE,
	headpunct=.,
	postheadspace=1em,
	spaceabove=15pt,spacebelow=10pt,
	postheadhook=\theorembookmark]%
{mainstyle}
\declaretheoremstyle[%
	headfont=\sffamily\bfseries,
	notefont=\normalfont\sffamily\scshape,
	bodyfont=\normalfont,
	headformat=\NUMBER\NAME\NOTE,
	headpunct=.,
	postheadspace=1em,
	spaceabove=15pt,spacebelow=10pt,
	postheadhook=\theorembookmark]%
{mainstyle_unnumbered}
\declaretheoremstyle[%
	headfont=\sffamily\bfseries,
	notefont=\normalfont\sffamily\scshape,
	bodyfont=\normalfont,
	headformat=swapnumber,
	headpunct=.,
	postheadspace=1em,
	spaceabove=15pt,spacebelow=10pt,
	postheadhook=\theorembookmark,
	qed=\qedsymbol]%
{mainstyleB}
\declaretheorem[name=Definition,parent=section,style=mainstyle]{definition}
\declaretheorem[name=Definition,numbered=no,style=mainstyle_unnumbered]{definition*}
\declaretheorem[name=Theorem,sharenumber=definition,style=mainstyle]{theorem}
\declaretheorem[name=Theorem,numbered=no,style=mainstyle_unnumbered]{theorem*}
\declaretheorem[name=Proposition,sharenumber=definition,style=mainstyle]{proposition}
\declaretheorem[name=Lemma,sharenumber=definition,style=mainstyle]{lemma}
\declaretheorem[name=Satz,sharenumber=definition,style=mainstyle]{satz}
\declaretheorem[name=Korollar,sharenumber=definition,style=mainstyle]{korollar}
\declaretheorem[name=Korollar,sharenumber=definition,style=mainstyleB]{korollarB}
\declaretheorem[name=Frage,numbered=no,style=mainstyle_unnumbered]{frage}
\declaretheorem[name=Erinnerung,sharenumber=definition,style=mainstyle]{erinnerungA}


% -- Beweise
\declaretheoremstyle[headfont=\bfseries\scshape,bodyfont=\normalfont,headpunct=:,postheadspace=1em,spacebelow=12pt,spaceabove=2pt,qed=\qedsymbol]{beweise}
\declaretheoremstyle[headfont=\sffamily\bfseries,bodyfont=\normalfont,headpunct=:,postheadspace=1em,spacebelow=10pt,spaceabove=10pt]{bemerkungen}
\declaretheorem[name=Beweis,numbered=no,style=beweise]{beweis}
\declaretheorem[name=Bemerkung,sharenumber=definition,style=mainstyle]{bemerkung}
\declaretheorem[name=Beispiel,sharenumber=definition,style=mainstyle]{beispiel}
\declaretheorem[name=Übung,numbered=no,style=bemerkungen]{uebung}
\declaretheorem[name=Erinnerung,numbered=no,style=bemerkungen]{erinnerung}


%--Skalarprodukt
\DeclarePairedDelimiterX\sprod[2]{\langle}{\rangle}{#1\,\delimsize\vert\,#2}
\DeclarePairedDelimiterX\skal[2]{\langle}{\rangle}{#1\,,\,#2}

%--Betrag, Gaußklammer
\DeclarePairedDelimiter{\abs}{\lvert}{\rvert}
\DeclarePairedDelimiter{\floor}{\lfloor}{\rfloor}
\DeclarePairedDelimiter{\ceil}{\lceil}{\rceil}

%--Norm
\DeclarePairedDelimiter\norm{\Vert}{\Vert}


%--Umklammern
\DeclarePairedDelimiter\enbrace{(}{)}
\DeclarePairedDelimiter\benbrace{[}{]}
\DeclarePairedDelimiter\lenbrace{<}{>}
\newcommand{\ssbrace}[1]{{\scriptscriptstyle\enbrace{#1}}}

%--Mengen
\newcommand\SetSymbol[1][]{\nonscript\:#1\vert\allowbreak\nonscript\:\mathopen{}}
\providecommand\given{} % to make it exist
\DeclarePairedDelimiterX\set[1]\{\}{\renewcommand\given{\SetSymbol[\delimsize]}#1}

%--Abbildungsdefinition
\newcommand{\mapdef}[5]{%
	\[
		\begin{array}{rcl}
			\textstyle #1 &\xrightarrow{\minwidthbox{#5}{2em}} & \textstyle #2 \\[0.5ex]
			\textstyle #3 &\xmapsto{\minwidthbox{\mbox{ }}{2em}} & \textstyle #4
		\end{array}
	\]
}

%--modifiziertes Stackrel
\newcommand{\StackText}[2]{\stackrel{\mbox{\scriptsize #1}}{#2}}

%--Befehl für eine Box mit Mindestweite
\DeclareRobustCommand{\minwidthbox}[2]{%
  \ifmmode
    \expandafter\mathmakebox
  \else
    \expandafter\makebox
  \fi
  [\ifdim#2<\width\width\else#2\fi]{#1}%
}

%--Befehl zum Einkreisen
\newcommand*\circled[1]{\tikzexternaldisable\tikz[baseline=(char.base)]{\node[shape=circle,draw,inner sep=2pt] (char) {#1};}\tikzexternalenable}


%--Differential
\newcommand{\diff}[2]{\ensuremath{\frac{{\partial #1}}{{\partial #2}} }}
\newcommand{\diffd}[2]{\ensuremath{\frac{\mathd #1}{\mathd #2} }}

%--Inhaltsverzeichnis
\usepackage[tocindentauto]{tocstyle}
\usetocstyle{KOMAlike}	
\shorthandon{"}
		
%--Punkte (als letztes laden)
\usepackage{ellipsis}

\newcommand{\fach}{Operatoralgebren}
\newcommand{\semester}{WiSe 2015}
\newcommand{\homepage}{https://wwwmath.uni-muenster.de/u/wilhelm.winter/wwinter/operatoralgebren.html}

\newcommand{\prof}{Prof.\ Dr.\ Wilhelm Winter}
% \publishers{\scalebox{14}{$K$}}
\providecommand{\verfasser}{Jannes Bantje}

%--Konfiguration von scrheadings
\setheadsepline{1pt}[\color{light_gray}]
\pagestyle{scrheadings}
%\fancyhf[H, F]{}
\clearscrheadfoot

\providecommand{\shortFach}{\fach}
\lehead{ \includegraphics[height=0.5 cm,keepaspectratio]{../!config/Bilder/Logo_WWU_Muenster_light_gray.pdf}}
\rehead{\rule{0cm}{0.5cm}\footnotesize \sffamily \color{light_gray} \verfasser -- Mitschrift \shortFach}
\lohead{\rule{0cm}{0.5cm} \footnotesize \sffamily \color{light_gray} Stand: \today \; \thistime[:]}
\rohead{\includegraphics[height=0.5 cm,keepaspectratio]{../!config/Bilder/fb10logo_gray.pdf}}


\ofoot[{ \color{dark_gray} \LARGE \sffamily \thepage}]{{ \color{dark_gray} \LARGE \sffamily \thepage}} %hier wir auch der plain Stil bearbeitet!
\automark{section}
\ifoot{ \color{dark_gray} \small \leftmark}

%--Metadaten
\author{\verfasser}
\titlehead{\includegraphics[height=1.5cm, keepaspectratio]{../!config/Bilder/Logo_WWU_Muenster.pdf}%
\hfill \includegraphics[height=1.3cm, keepaspectratio]{../!config/Bilder/fb10logo.pdf}}
\title{Skript \fach}
\subtitle{Mitschrift der Vorlesung  \enquote{\fach} von \prof}

\publishers{
Aktuelle Version verfügbar bei: \bigskip\\
\normalsize
\begin{minipage}[t]{0.45\textwidth}
	\centering{
	\qrcode[height=3.5cm, version=4]{https://github.com/JaMeZ-B/latex-wwu}
	\medskip\\
	\includegraphics[height=0.4cm, keepaspectratio]{../!config/Bilder/github_octo.pdf} 
	\includegraphics[height=0.4cm, keepaspectratio]{../!config/Bilder/GitHub_Logo.pdf} (inklusive Sourcecode)\\
	{\small\url{https://github.com/JaMeZ-B/latex-wwu}}}
\end{minipage}
\hfill
\begin{minipage}[t]{0.45\textwidth}
	\centering{
	\qrcode[height=3.5cm, version=4]{btsync://B6WH2DISQ5QVYIRYIEZSF4ZR2IDVKPN3I?n=latex_share}
	\medskip\\
	\includegraphics[height=0.4cm, keepaspectratio]{../!config/Bilder/bt_sync_logo.pdf} 
	{\large \textbf{Bittorrent} Sync} \\
	\texttt{B6WH2DISQ5QVYIRYIEZSF4ZR2IDVKPN3I}
	}
\end{minipage}
}


\begin{document}
\pagenumbering{Roman}
\maketitle
\begin{abstract}
\section*{Vorwort --- Mitarbeit am Skript}
Dieses Dokument ist eine Mitschrift aus der Vorlesung \enquote{\fach, \semester}, gelesen von \prof. Der Inhalt entspricht weitestgehend dem Tafelanschrieb. Für die
Korrektheit des Inhalts übernehme ich keinerlei Garantie! Für Bemerkungen und Korrekturen -- und seien es nur Rechtschreibfehler -- bin ich sehr dankbar. 
Korrekturen lassen sich prinzipiell auf drei Wegen einreichen: 
\begin{itemize}
	\item Persönliches Ansprechen in der Uni, Mails an \hrefsymmail{mailto:\mail}{\mail} (gerne auch mit annotieren PDFs) 
	\item \emph{Direktes} Mitarbeiten am Skript: Den Quellcode poste ich auf GitHub (siehe oben), also stehen vielfältige Möglichkeiten der Zusammenarbeit zur Verfügung:
	Zum Beispiel durch Kommentare am Code über die Website und die Kombination Fork + Pull Request. Wer sich verdient macht oder ein Skript zu einer Vorlesung, die 
	ich nicht besuche, beisteuern will, dem gewähre ich gerne auch Schreibzugriff.
	
	Beachten sollte man dabei, dass dazu ein Account bei \url{github.com} notwendig ist, der allerdings ohne Angabe von persönlichen Daten angelegt werden kann. 
	Wer bei GitHub (bzw. dem zugrunde liegenden Open-Source-Programm \enquote{\texttt{git}}) -- verständlicherweise -- Hilfe beim Einstieg braucht, dem helfe ich gerne 
	weiter. Es gibt aber auch zahlreiche empfehlenswerte Tutorials im Internet.\footnote{zB. \url{https://try.github.io/levels/1/challenges/1}, ist auf Englisch, aber dafür 
	interaktives LearningByDoing}
	\item \emph{Indirektes} Mitarbeiten: \TeX-Dateien per Mail verschicken. 
	
	Dies ist nur dann sinnvoll, wenn man einen ganzen Abschnitt ändern möchte (zB. einen alternativen Beweis geben), da ich die Änderungen dann per Hand einbauen muss! Ich freue mich aber auch über solche Beiträge!
\end{itemize}
\section*{Vorlesungshomepage}
{\centering 
\begin{minipage}[c][][c]{\textwidth}
	\centering \qrcode[height=3cm]{\homepage} \medskip\\
	\footnotesize \url{\homepage}
\end{minipage}
\par}

\section*{Anmerkung}
Innerhalb dieser Mitschrift wird man öfter den Ausdruck \enquote{Warum?} finden. Dies sind vom Dozenten bewusst weggelassene Details, die zu verstärktem Mitdenken beim Lesen 
animieren sollen. Oftmals sind dies schon aus vorherigen Semestern bekannte Sachverhalte. Nur an wenigen Stellen habe ich die fehlenden Details hinzugefügt.
\end{abstract}

\tableofcontents
\cleardoubleoddemptypage

\pagenumbering{arabic}
\setcounter{page}{1}
\setcounter{footnote}{0}

\section{Spektraltheorie in Banachalgebren} % (fold)
\label{sec:1}

\begin{definition}[{name=[Banachalgebra]}]
	Eine \Index{Banachalgebra} ist eine vollständige normierte Algebra. Für die Multiplikation gilt also $\norm{a b} \le \norm{a} \cdot \norm{b}$
	Wir betrachten in dieser Vorlesung fast ausschließlich Algebren über $\mathds{C}$. Eine Banachalgebra heißt \Index{unital}, falls sie ein Einselement $1$ mit $\norm{1} =1$ 
	enthält. 
\end{definition}

\begin{definition}[{name=[Menge der invertierbaren Elemente und Spektrum]}]
	Für eine unitale $\mathds{C}$-Algebra $A$ setzen wir 
	\begin{align}
		\Inv(A) &:= \set{a \in A  \text{ invertierbar} } \\
		\sigma(a) &:=\sigma_A(a) := \set[\big]{\lambda \in \mathds{C} \given \lambda \cdot 1 - a \notin \Inv(A)}
	\end{align}
	$\sigma(a)$ heißt \Index{Spektrum} von $a$. \todo{Revision}
\end{definition}

\begin{bemerkung}
	Sei $A$ wie oben und $a,b \in A$. Dann gilt
	\[
		1 -ab \in \Inv(A) \iff 1 -ba \in \Inv(A)
	\]
\end{bemerkung}
\begin{beweis}
	Durch simplex Nachrechnen verifiziert man, dass $(1-ba)^{-1} = 1+ b(1- ab)^{-1} a$ gilt.
\end{beweis}

\begin{satz}[label=satz14_poly]
	Sei $A$ eine unitale $\mathds{C}$-Algebra und $a \in A$ mit $\sigma(a)\not= \emptyset$. Sei weiter $p \in \mathds{C}[Z]$ ein Polynom. Dann gilt 
	\[
		\sigma\enbrace[\big]{p(a) }= p \enbrace[\big]{\sigma(a)}
	\]
\end{satz}
\begin{beweis}
	Die Aussage ist trivial, falls $p$ konstant ist. Sei also $p$ nicht konstant. Zu $\mu \in \mathds{C}$ existieren $\lambda_0, \ldots , \lambda_n \in \mathds{C}$ mit 
	$\lambda_0 \not= 0$ und $p(z) -\mu = \lambda_0 \cdot (z- \lambda_1) \cdot  \ldots \cdot (z- \lambda_n)$, also
	\[
		p(a) - \mu = \lambda_0 \cdot (a-\lambda_1) \cdot  \ldots \cdot (a- \lambda_n)
	\]
	Damit gilt: $p(a) -\mu$ ist genau dann invertierbar, wenn $a -\lambda_1, \ldots , a-\lambda_n $ invertierbar sind. Also 
	\begin{align}
		\mu \in \sigma \enbrace[\big]{p(a)} &\iff \lambda_i \in \sigma(a) \text{ ist für ein } i \in \set{1,\ldots ,n} \\
		&\iff p(\lambda) -\mu =0 \text{ für ein } \lambda \in \sigma(a) \\
		&\iff p(\lambda)=\mu \text{ für ein } \lambda \in \sigma(a)   
	\end{align}
	Somit folgt dann $\sigma \enbrace*{p(a)} = p \enbrace*{\sigma(a)}$.
\end{beweis}

\begin{satz}[label=satz15_neumann]
	Sei $A$ eine unitale Banachalgebra und $a \in A$. Falls $\norm{a}<1$, so ist $1-a$ invertierbar und 
	\[
		(1-a)^{-1} = \sum_{n=0}^{\infty} a^n \tag*{(\Index{Neumannsche Reihe})}
	\]
\end{satz}
\begin{beweis}
	Die Folge der Partialsummen $\enbrace[\big]{\sum_{n=0}^{k} a^n}_{k \in \mathds{N}} \subset A$ ist eine Cauchyfolge (warum?), sodass die unendliche Summe existiert. Es gilt dann
	\begin{align}
		(1-a) \enbrace*{\sum_{n=0}^{\infty} a^n} \xleftarrow{k \to \infty}(1-a) \enbrace*{\sum_{n=0}^{k} a^n} = 1-a^{k+1} \xrightarrow{k \to \infty} 1
	\end{align}
	Die Summe ist also ein Rechtsinverses. Genauso folgt, dass es auch ein Linksinverses ist.
\end{beweis}

\begin{korollar}[{name=[Invertierbare Elemente bilden offene Menge]},label=inv_offen]
	Für eine unitale Banachalgebra $A$ gilt $\Inv(A) \subset A$ offen.
\end{korollar}
\begin{beweis}
	Falls $a \in \Inv(A)$ und $\norm{b-a} < \norm{a^{-1}}^{-1}$, so gilt $\norm{b a^{-1} -1} \le \norm{b-a} \norm{a^{-1}}<1$, also ist 
	$-\enbrace*{1- \enbrace*{ba^{-1} -1}} = ba^{-1}$ invertierbar nach \autoref{satz15_neumann} und damit auch $b \in \Inv(A)$, denn $b$ ist ein Linksinverses für 
	$a^{-1}(ba^{-1})^{-1} \in \Inv(A)$.
\end{beweis}

\begin{definition}[{name=[Spektralradius]}]
	Sei $A$ eine unitale Banachalgebra. Für $a \in A$ definieren wir den \Index{Spektralradius} durch
	\[
		r(a) := \sup_{\lambda \in \sigma(a)} \abs{\lambda} 
	\] 
\end{definition}

\begin{proposition}[{name=[Spektrum ist kompakt und Spektralradius durch Norm beschränkt]},label=spec_kompakt]
	Seien $A,a$ wie oben. Dann ist $\sigma(a)$ kompakt und $r(a) \le \norm{a}$.
\end{proposition}
\begin{beweis}
	Angenommen $\abs{\lambda} > \norm{a}$, dann gilt $\norm{\lambda^{-1} \cdot a} < 1$, also ist $1- \lambda^{-1} \cdot a$ invertierbar nach \autoref{satz15_neumann} und 
	$\lambda -a \in \Inv(A)$. Damit ist $\lambda \notin \sigma(a)$. Damit folgt, dass $r(a) \le \norm{a}$ ist. 	
	Es ist $\mathds{C} \setminus \sigma(a) \subset \mathds{C}$ offen, da $\Inv(A) \subset A$ offen ist. Damit muss $\sigma(a)$ kompakt sein.
\end{beweis}

\begin{satz}[{name=[In unitalen Banachalgebren sind die Spektren nie leer]},label=spec_nichtleer]
	Sei $A$ eine unitale Banachalgebra. Für jedes $a \in A$ ist $\sigma(a)$ nichtleer. 
\end{satz}
\begin{beweis}
	Falls $\abs{\lambda}> \norm{a}$, so existiert $\enbrace*{1- \lambda^{-1} a}^{-1}$, also $(\lambda-a)^{-1}$ und es gilt
	\begin{align}
		(\lambda-a)^{-1} = \lambda^{-1} \enbrace*{1- \lambda^{-1} a}^{-1} = \lambda^{-1} \enbrace*{\sum_{n=0}^{\infty} \lambda^{-n}a^n}
	\end{align}
	Es folgt 
	\[
		\norm{(\lambda-a)^{-1}} \le \abs*{\lambda^{-1}} \cdot \Underbrace{\sum_{n=0}^{\infty} \enbrace*{\abs*{\lambda^{-1}} \cdot \norm{a}}^n}{\le 2 
		\text{ , falls } \abs{\lambda}> 2 \norm{a}} \xrightarrow{\abs{\lambda} \to \infty} 0 \label{19:eq:1} \tag{$*$}
	\]
	Sei nun $\lambda_0 \in \mathds{C}\setminus \sigma(a)$ und $\abs{\lambda-\lambda_0}< \norm{(a- \lambda_0)^{-1}}^{-1}$. Dann gilt 
	\[
		\norm{(\lambda-\lambda_0)(a - \lambda_0)^{-1}} \le \norm{\lambda-\lambda_0} \norm{(a-\lambda_0)^{-1}} <1   
	\]
	also $\enbrace*{1- (\lambda-\lambda_0)(a-\lambda_0)^{-1}}^{-1} = \sum_{n=0}^{\infty} (\lambda-\lambda_0)^n (a-\lambda_0)^n$ und 
	\begin{align}
		-\sum_{n=0}^{\infty} (\lambda-\lambda_0)^n(a-\lambda_0)^{-(n+1)} &= (a- \lambda_0)^{-1} \enbrace*{(\lambda-\lambda_0)(a-\lambda_0)^{-1}-1}^{-1} \\
		&= \enbrace[\Big]{(\lambda-\lambda_0) \Underbrace{(a-\lambda_0)(a-\lambda_0)^{-1}}{=1}-(a- \lambda_0)}^{-1} = (\lambda-a)^{-1}
	\end{align}
	Für ein stetiges Funktional $\varphi \in A^*$ erhalten wir 
	\[
		\varphi\enbrace*{(\lambda-a)^{-1}} = \sum_{n=0}^{\infty} - \varphi \enbrace*{(a-\lambda_0)^{-(n+1)}} \cdot (\lambda-\lambda_0)^n
	\]
	falls $\abs{\lambda-\lambda_0} < \norm{(a- \lambda_0)^{-1}}^{-1}$. Daher ist $\lambda \mapsto \varphi \enbrace*{(\lambda-a)^{-1}}$ \emph{holomorph} auf 
	$\mathds{C}\setminus \sigma(a)$. Falls $\sigma(a)= \emptyset$, so ist $\lambda \mapsto \varphi \enbrace*{(\lambda-a)^{-1}}$ \emph{ganz}. Wegen \eqref{19:eq:1} und dem Satz
	von Liouville\footnote{siehe \url{https://de.wikipedia.org/wiki/Satz_von_Liouville_\%28Funktionentheorie\%29}} ist dann $\varphi \enbrace*{(\lambda-a)^{-1}}=0$ für alle 
	$\lambda\in \mathds{C}$. Aber $\varphi \in A^*$ war beliebig, und wegen Hahn-Banach gilt $(\lambda-a)^{-1}=0$. Dies ist ein Widerspruch.
\end{beweis}

\begin{satz}
	Sei $A$ eine unitale Banachalgebra. Für $a \in A$ gilt dann 
	\[
		r(a)=\lim_{n \to \infty}\norm{a^n}^{\frac{1}{n}}
	\]
\end{satz}
\begin{beweis}
	Sei $\varphi \in A^*$ ein stetiges Funktional. Wir haben im Beweis von \autoref{spec_nichtleer} gesehen, dass die Zuordnung 
	$\lambda \mapsto \varphi \enbrace*{(\lambda-a)^{-1}}$ analytisch ist für $\lambda \notin \sigma(a)$. Insbesondere gilt dies für $\abs{\lambda}> r(a)$. Dann ist auch
	$\lambda \mapsto \varphi \enbrace*{(\lambda^{-1}-a)^{-1}}$ analytisch für $0< \abs{\lambda}< \sfrac{1}{r(a)}$, ebenso 
	$\lambda \mapsto \varphi \enbrace*{(1- \lambda \cdot a)^{-1}}$. Die Singularität bei $0$ ist hebbar, also definiert $f(\lambda) := \varphi \enbrace*{(1-\lambda \cdot a)^{-1}}$
	eine analytische Funktion für $\abs{\lambda} < \sfrac{1}{r(a)}$. Dann existiert eine Folge $(\alpha_n)_{n \in N} \subset \mathds{C}$ mit 
	$f(\lambda) = \sum_{n=0}^{\infty} \alpha_n \cdot \lambda^n$ für $\abs{\lambda}<r(a)$. Die Konvergenz ist absolut und gleichmäßig, also gilt insbesondere 
	$\abs{\alpha_n \cdot \lambda^n} \xrightarrow{n \to \infty}0$. Falls $\abs{\lambda} < \sfrac{1}{\norm{a}} \le \sfrac{1}{r(a)}$ 
	ist, so gilt
	\[
		f(\lambda) = \varphi \enbrace*{(1- \lambda \cdot a)^{-1}} \stackrel{\mbox{\scriptsize\ref{satz15_neumann}}}{=} \varphi \enbrace*{\sum_{n=0}^{\infty} (\lambda a)^n} = 
		\sum_{n=0}^{\infty} \varphi \enbrace[\big]{(\lambda a)^n} = \sum_{n=0}^{\infty} \varphi(a^n) \cdot \lambda^n
	\]
	also $\alpha_n = \varphi(a^n)$ für $n \in N$. Es gilt $\varphi(a^n) \cdot \lambda^n \xrightarrow{n \to \infty} 0$. Aber $\varphi \in A^*$ war beliebig, also ist für jedes 
	$\varphi \in A^*$ die Folge $\enbrace*{\varphi(\lambda^n a^n)}_\mathds{N} = \enbrace*{\varphi(a^n) \cdot \lambda^n}_{n \in \mathds{N}}$ beschränkt für 
	$\abs{\lambda}<\sfrac{1}{r(a)}$.

	Nach dem Prinzip der gleichmäßigen Beschränktheit ist dann $(\lambda^n a^n)_{n \in \mathds{N}}$ beschränkt für jedes $\lambda$ mit $\abs{\lambda}< \sfrac{1}{r(a)}$. Damit 
	existiert für $0 < \abs{\lambda}< \sfrac{1}{r(a)}$ ein $M \ge 0$ mit $\norm{\lambda^n a^n} \le M$, $n \in \mathds{N}$, also 
	\[
		\norm{a^n}^{1/n}\le\frac{M^{\sfrac{1}{n}}}{\abs{\lambda}} \xrightarrow{n \to \infty} \frac{1}{\abs{\lambda}} 
	\]
	Es folgt $\limsup_n \norm{a^n}^{\sfrac{1}{n}} \le r(a)$. Falls $\lambda \in \sigma(a)$, so gilt nach \autoref{satz14_poly} $\lambda^n \in \sigma(a^n)$ und wegen 
	\autoref{spec_kompakt} $\abs{\lambda^n} \le \norm{a^n}$, also $\abs{\lambda} \le \norm{a^n}^{\sfrac{1}{n}}$ und es folgt insgesamt
	\begin{align}
		r(a) \le \inf_n \norm{a^n}^{\sfrac{1}{n}} \le \liminf_n \norm{a^n}^{\sfrac{1}{n}} \le \limsup_n \norm{a^n}^{\sfrac{1}{n}} \le r(a)
	\end{align}
	und damit $\lim_n \norm{a^n}^{\sfrac{1}{n}} = r(a)$.
\end{beweis}

\begin{proposition}[label=prop1.11_einfach]
	Sei $A$ eine abelsche unitale Banachalgebra. Falls $A$ einfach\footnote{besitzt keine nichttrivialen Ideale} ist, so ist $A \cong \mathds{C}$. 
\end{proposition}
\begin{beweis}
	Sei $a \in A$, $\alpha \in \sigma(a)$. Dann ist $J := \overline{(a-\alpha \cdot 1)\cdot A} \lhd A$ ein abgeschlossenes Ideal. Für jedes $b \in A$ ist $(a-\alpha)b$ nicht 
	invertierbar. Damit folgt dann $\norm{(a-\alpha)b - 1}\ge 1$ für jedes $b \in A$. Damit ist auch $\dist\enbrace*{(a-\alpha \cdot 1)\cdot A,1}\ge 1$. Also ist $1 \notin J$ und 
	damit $j \not= A$. Da $A$ einfach ist, folgt $J=\set{0}$, also muss $a-\alpha \cdot 1=0$ und weiter $a \in \mathds{C} \cdot 1$ gelten. Damit ist 
	$A = \mathds{C} \cdot 1 \cong \mathds{C}$.
\end{beweis}

\begin{definition}[{name=[Charakter]}]
	Sei $A$ eine abelsche unitale Banachalgebra. Eine lineare, multiplikative, nicht verschwindende Abbildung $\chi \colon A \to \mathds{C}$ heißt \Index{Charakter} auf $A$.  
	Es gilt dann $\chi(1)=1$. Wir setzen $\Omega(A) := \set{\chi \given \chi \text{ Charakter auf }A}$ (\Index{Charakterspektrum} oder einfach nur \Index{Spektrum} von $A$).
	
	Die \Index{Geldfanddarstellung} ist definiert als 
	$a \mapsto \hat{a} \colon \Omega(A) \to \mathds{C}, \chi \mapsto \chi(a)$.
\end{definition}

\begin{satz}
	Sei $A$ eine abelsche unitale Banachalgebra. Dann gilt
	\begin{enumerate}[(i)]
		\item Jeder Charakter auf $A$ ist stetig mit Norm $1$.
		\item Die Abbildung $\chi \mapsto \ker \chi$ definiert eine Bijektion zwischen 
		\[
			\Omega(A) \to J_{\max} := \set[\big]{J \given J \subsetneq A \text{ ist maximales Ideal} }
		\]
	\end{enumerate}
\end{satz}
\begin{beweis}\leavevmode
	\begin{enumerate}[(i)]
		\item Sei $a \in A$ mit $\norm{a}<1$ und $\chi(a)=1$. Setze $b= \sum_{n=1}^{\infty} a^n$, dann ist $1 + b = \sum_{n=0}^{\infty} a^n = (1-a)^{-1}$. Also ist
		$1 = (1+b)(1-a) = 1 + b -a -ab$ und damit $ab+a=b$. Es folgt $\chi(a)\chi(b) + \chi(a)= \chi(b)$. Dies ist ein Widerspruch. Also muss $\norm{\chi}\le 1$ sein. Mit 
		$\chi(1)=1$ folgt $\norm{\chi}=1$.
		\item Für jedes $\chi \in \Omega(A)$ ist $\ker \chi \lhd A$ ein Ideal. Weiter ist der Kern abgeschlossen in $A$, da $\chi$ stetig ist. Wir zeigen, dass $\ker \chi$ maximal 
		ist. $\ker \chi \not= A$, denn $\chi \not= 0$. Falls $c \in A \setminus \ker \chi$, so gilt 
		\(
			\chi(c)^{-1} \cdot c-1 \in \ker \chi
		\).
		Also ist $\ker \chi + \mathds{C} \cdot c = \ker \chi + \mathds{C} \cdot 1=A$, da $a- \chi(a) \cdot 1 \in \ker \chi$ und damit 
		$a= \enbrace*{a- \chi(a) \cdot 1}+ \chi(a) \cdot 1 \in \ker \chi + \mathds{C} \cdot 1$. Damit ist der Kern von $\chi$ maximal.
		
		Zur Injektivität: Es seien $\chi_1, \chi_2 \in \Omega(A)$ mit $\ker \chi_1 = \ker \chi_2$. Dann gilt für $a \in A$
		\begin{align}
			\chi_1(a) - \chi_2(a) = \chi_1 \enbrace[\Big]{\Underbrace{a - \chi_2(a) \cdot 1}{\in \ker \chi_2=\ker \chi_1}} = 0
		\end{align}
		Also ist $\chi_1=\chi_2$. Für die Surjektivität betrachten wir ein maximales, nichttriviales Ideal $J \subsetneq A$. Dann ist $\dist \enbrace*{J,1}=1$, denn 
		\todo{Einsen auseinanderhalten} 
		$\mathring{B}_A(1,1) \subset \Inv(A)$ nach \autoref{satz15_neumann}. Dann ist $1 \notin \overline{J} \lhd A$. Da $J$ maximal ist, folgt $J=\overline{J}$. Damit ist 
		$\sfrac{A}{J}$ eine einfache, abelsche, unitale Banachalgebra. Mit \autoref{prop1.11_einfach} folgt $\sfrac{A}{J} \cong \mathds{C}$. Damit ist 
		$\chi \colon A \to \sfrac{A}{J} \cong \mathds{C}$ ein Charakter mit $\ker \chi= J$. \qedhere
	\end{enumerate}
\end{beweis}
% section 1 (end)






\cleardoubleoddemptypage
\pagenumbering{Alph}
\setcounter{page}{1}

\printindex
\listoffigures
\todototoc
\listoftodos[To-do's und andere Baustellen]
\end{document}